%%%%%%%%%%  scan idx 233  |  folio 220  %%%%%%%%%%
\subsection*{BIBLIOGRAPHIE}
\begin{enumerate}[label={[\arabic*]}]
\item
Artin, M., Grothendieck, A., et Verdier, J-L., Cohomologie étale des schémas, à paraître dans North Holl. Pub. Comp., (cité SGA 4).
\item Atiyah, M. et Segal, G., papiers secrets.
\item Cartan, H. Funzioni e varietà complesse, Faisceaux analytiques cohérents, Cours au Centro Internazionale Matematico Estivo, 1964, Roma.
\item Cartan, H., et Schwartz, L., Théorème d'Atiyah--Singer, Séminaire 1963-64, Institut Henri Poincaré, Paris.
\item Frisch, J., Fonctions analytiques sur un ensemble semi-analytique, note aux CRAS Paris, t. 260, p. 2974-2976 (1965).
\item Godement, R., Théorie des Faisceaux, Hermann (1958), (cité T.F).
\item Grothendieck, A., Cohomologie locale des faisceaux cohérents et théorèmes de Lefschetz locaux et globaux, Séminaire de Géométrie Algébrique de l'IHES 1962, à paraître dans North Holl. Pub. Comp., (cité SGA 2).
\item Grothendieck, A., Fondements de la Géométrie Analytique, Séminaire Cartan 1959-60, Institut Henri Poincaré, Paris.
\item Hartshorne, R., Residues and Duality, Lecture Notes, Springer n\textsuperscript{o}~20, cité [H].
\item Houzel, C., Géométrie Analytique Locale, Séminaire Cartan 1959-60, Institut Henri Poincaré, Paris.
\item Illusie, L., Complexes de fibrés et définitions du foncteur $K$, Séminaire Shih Weishu, IHES 1964-65.
\item Illusie, L., Complexes quasi-acycliques directs de fibrés banachiques, note aux CRAS Paris, t. 260, p. 6499-6502.
\item Kiehl, R., Theorem A und Theorem B in der nichtarchimedischen Funktionentheorie, Invent. math. 2, p. 256-273 (1967).
